Le tableau ci-dessous indique \emph{la solubilités} de divers gaz dans l'eau à
$\qty{20}{\degreeCelsius}$ sous une pression de $\qty{1e5}{\Pa}$. Elle est
numériquement égale au volume maximal de gaz (en \unit{\L}) que l'on peut
dissoudre par litre d'eau. On obtient ainsi une solution saturée.

\begin{minipage}{\linewidth}
    \begin{table}[H]
    \centering
    \setlength{\arrayrulewidth}{0.8pt}
    \renewcommand{\arraystretch}{1.5}
    \begin{tabularx}{\textwidth}{|p{6.5cm}|C|C|C|}
        \hline
        Gaz & \ce{CO2} & \ce{HCl} & \ce{NH3} \\
        \hline
        $s$ (en \unit{\L} de gaz dissous par \unit{\L} d'eau) & 0,900 & 445 &
        800 \\
        \hline
    \end{tabularx}
\end{table}
\end{minipage}

\begin{enumerate}
    \item Calculer les concentrations molaires des solutions saturées en ces
          gaz.
    \item À partir de solutions saturées en \ce{HCl} et \ce{NH3}, on désire
          préparer $\qty{1.00}{\L}$ de solutions de concentration
          $\qty{1.00}{\mol\per\L}$. Quels volumes faut-il en prélever~?
\end{enumerate}

